\documentclass[11pt]{article}
\usepackage[T1]{fontenc}
\usepackage{lmodern,amsmath,amssymb,microtype}
\usepackage[margin=1.02in]{geometry}
\usepackage[colorlinks=true,urlcolor=blue,citecolor=blue]{hyperref}
\setlength{\emergencystretch}{2em}
\hypersetup{pdfauthor={Wentao Li},pdftitle={Rigidity need not descend from a power set}}
\title{Rigidity need not descend from a power set}
\author{Wentao Li\\[2pt]\small Prepared with assistance from OpenAI Codex}
\date{November 2026}
\begin{document}
\maketitle
\begin{abstract}
Relative to the consistency of ZF, there is an amorphous set $A$ with no rigid binary relation whose power set carries a rigid binary relation. An internally rigid ternary structure on $A$ supplies the additional information. This gives a negative answer to Howard and Tachtsis, Problem 7.
\end{abstract}
\section{A permutation model}
Start with a ground universe with choice and a countably infinite set $A$ of atoms. Externally identify $A$ with a countably infinite-dimensional vector space over $\mathbb F_2$. Use the group $G=\mathrm{GL}(A)$ and the normal filter generated by pointwise stabilizers of finite atom sets. Work in the hereditarily symmetric model $M$. Addition and zero belong to $M$, being invariant under $G$, with hereditarily symmetric constituents. Every internal object has a finite support.
If $B\subseteq A$ is internal with support $E$, the stabilizer of $E$ acts transitively on the complement of the finite subspace $U=\langle E\rangle$. Indeed two vectors outside $U$ extend a basis of $U$ by a single new independent vector in the same way, and the partial linear isomorphism extends externally to the whole space. Hence $B$ is finite or cofinite. The atoms are amorphous. In particular none of their binary relations is rigid, by the following ZF argument. Complement each cofinite row and record a unary marker for it. Finite row sizes are bounded, since an infinite image in $\omega$ would split the domain into two infinite sets. Only finitely many columns can be cofinite, by intersecting any collection of them of size larger than the row bound. Remove their indices and mark incidences with them as unary colors. The remaining underlying graph is locally finite, and each component is finite since otherwise its odd and even distance levels give disjoint infinite subsets. The component sizes are bounded by the same finite-image argument. Two components have isomorphic finite colored structures, so interchanging them gives a nonidentity automorphism. This uses only sets within $M$.
\section{Rigidity of addition inside the model}
Let $f$ be an internal linear automorphism, with finite support $E$. It commutes with every $h\in G$ fixing $E$ pointwise. For $v\notin U$, the vectors fixed by the stabilizer of $U\cup\{v\}$ are exactly $U+\mathbb F_2v$, since the group acts transitively on its infinite complement. Therefore $f(v)=c+\epsilon v$ for $c\in U$ and $\epsilon\in\{0,1\}$. By transitivity on $A\setminus U$ and commutation, the same $c,\epsilon$ work at every outside vector. Injectivity forces $\epsilon=1$. Take $v,w$ independent modulo $U$. Then $f(v+w)=v+w+c=f(v)+f(w)=v+w$, so $c=0$. For $u\in U$, apply linearity to $v$ and $v+u$ to see that $f(u)=u$. The sole internal linear automorphism is the identity.
\section{A relation on the power set}
Use $\mathcal P^M(A)$, the model's own power set. Let $\mathcal T$ consist of the triples $\{u,v,u+v\}$ for distinct nonzero $u,v$. On $\mathcal P^M(A)$ put a directed edge from $\{a\}$ to each non-singleton $S$ containing $a$, and a loop precisely at each member of $\mathcal T$. This definition gives an internal relation. Singletons are exactly the sources of off-diagonal edges, since $A$ itself is always available as a target. A non-singleton is uniquely determined among the non-singletons by its incoming singleton neighbors, including the case of the empty set. Every automorphism therefore comes from an internal permutation of $A$ preserving the triples. Zero is the only point of $A$ in no triple. Distinct nonzero vectors have sum the unique third point of their triple. The cases of equal vectors or zero are determined as well, so the permutation is additive and hence the identity. This proves rigidity.
\paragraph{Pure ZF.} The Jech--Sochor embedding theorem transfers the iterated power-set structure over the atoms through any prescribed ordinal height. Apply it high enough for subsets of $A$, binary relations, bijections on $\mathcal P(A)$, and the assertions of finiteness for subsets of $A$ (height $\omega+\omega$ suffices, using relative superstructures). It preserves the full powersets at the relevant stages. Consequently the image of $A$ has exactly the stated three properties in a model of pure ZF. This is a relative-consistency assertion; in ZFC every set admits a rigid well-order.
\paragraph{Sources.} P. Howard and E. Tachtsis, \href{https://doi.org/10.4064/fm960-12-2015}{Fund. Math. 232 (2016)}, Theorem 4.10 gives the upward implication and Problem 7 asks about descent. Its singleton coding motivates the construction here. J. K. Truss, \href{https://www.sciencedirect.com/science/article/pii/016800729400024W}{\emph{The structure of amorphous sets}} (1995), treats finite-field geometries. Y. Hu, R. Mao and G. Shen, \url{https://arxiv.org/abs/2510.13508}, Section 2, use exactly the vector-space model above and prove its amorphousness for a different power-set question. The new application uses the internal centralizer and the triple incidence relation. For Jech--Sochor see T. Jech, \emph{The Axiom of Choice}, Theorem 6.1, or L. Halbeisen, \href{https://people.math.ethz.ch/~halorenz/publications/pdf/cst.pdf}{\emph{Combinatorial Set Theory}}, chapter on permutation models. Independent proof and priority review is pending.
\end{document}
